% ==========================================================================
% INFO 4617 Web Data Science -- Week 7 Friday activity: Review a pull request
% Build: `make` in handouts/, or `latexmk -pdf review-a-pull-request.tex`
% here. The figures are annotated from real captures (img/IMAGES.md). A
% screen that needs a signed-in browser is a hand capture: \handcapture
% prints it once img/NAME.png exists, and leaves it out until then.
% ==========================================================================
\makeatletter\def\input@path{{../common/}{handouts/common/}}\makeatother
\documentclass{handout}
\usepackage{amssymb}   % \checkmark

\title[INFO 4617 · Week 7 · Friday activities · Review a pull request]{Review a pull request}
\subtitle{Week 7 Friday activity · medium · about 15 minutes a pull request}
\author{Brian C. Keegan, Ph.D.}
\institute{Department of Information Science, University of Colorado Boulder}
\date{October 2, 2026}

% A numbered figure: [width]{path from the repo root}{alt text}{caption}
\newcounter{handoutfig}
\newcommand{\activityfigure}[4][\linewidth]{%
  \par\smallskip\refstepcounter{handoutfig}%
  {\centering
   \BeginAccSupp{method=pdfstringdef,Alt={#3}}%
   \includegraphics[width=#1]{#2}%
   \EndAccSupp{}\par}%
  \smallskip{\small\raggedright\textbf{Figure \thehandoutfig.} #4\par}}
% A hand capture: {file in handouts/week-07/img/}{alt text}{caption}
\newcommand{\handcapture}[3]{%
  \IfFileExists{handouts/week-07/img/#1}%
    {\activityfigure{handouts/week-07/img/#1}{#2}{#3}}%
    {\IfFileExists{../../handouts/week-07/img/#1}%
      {\activityfigure{handouts/week-07/img/#1}{#2}{#3}}{}}}

\setlist[enumerate]{leftmargin=*,itemsep=1.5pt,topsep=2pt}
\setlist[itemize]{leftmargin=*,itemsep=1pt,topsep=2pt}

\begin{document}
\maketitle

\begin{frame}{What a review is}
  \begin{columns}[T]
    \begin{column}{0.57\textwidth}
      A review reads a classmate's proposed change to the book. It says
      whether the change should go in, and what must change first. You read
      and you comment. You do not run code, and you do not change the pull
      request.

      Your chapter's board in week 7's slides lists the pull requests to
      review. A line like \texttt{\#88 and \#93} means that two pull
      requests make the same change. Read them together, and say which one
      to keep.

      Work in pairs. One partner works in the browser, and the other reads
      this handout aloud.
    \end{column}
    \begin{column}{0.39\textwidth}
      \begin{block}{The three Friday activities}
        \small
        \begin{tabular}{@{}ll@{}}
          Triage an issue & easy\\
          \textbf{Review a pull request} & \textbf{medium}\\
          Issue to pull request & hard
        \end{tabular}
      \end{block}
      \begin{alertblock}{Do not}
        \small
        \begin{itemize}
          \item click \textbf{Merge pull request}, even if GitHub shows it;
          \item change someone else's pull request: comment on it;
          \item close the pull request.
        \end{itemize}
        The instructor merges, closes, and resolves conflicts.
      \end{alertblock}
    \end{column}
  \end{columns}
\end{frame}

\begin{frame}{1 · Read what it proposes}
  \begin{enumerate}
    \item Sign in to GitHub.
    \item Open the pull request. Its address is
      \url{https://github.com/cuinfoscience/Web-Data-Science-Book/pull/}
      followed by its number, for example \texttt{.../pull/104}.
    \item On the \textbf{Conversation} tab, read the description. The
      template asks for \textbf{Closes \#N} (the issue it fixes),
      \textbf{Location}, \textbf{Problem}, \textbf{Why}, \textbf{Change},
      \textbf{AI assistance}, and \textbf{Checks}.
    \item If it names an issue, open the issue. Does the pull request do
      what the issue asks for?
    \item Read the comments. If another pair has reviewed it, read their
      review. Add only what is new.
  \end{enumerate}
\end{frame}

\begin{frame}{2 · Look at the checks}
  Scroll to the bottom of the \textbf{Conversation} tab.

  \smallskip
  {\small
  \begin{tabular}{@{}>{\raggedright\arraybackslash}p{0.33\linewidth}%
                  >{\raggedright\arraybackslash}p{0.63\linewidth}@{}}
    \toprule
    \textbf{You see} & \textbf{What it means}\\
    \midrule
    \textbf{Render} with a green \checkmark &
      The book still builds with this change.\\[3pt]
    \textbf{Render} with a red \texttimes &
      The book does not build. This is a must-fix: say so in your review.\\[3pt]
    \textbf{Notebook sync} with a red \texttimes &
      Normal for a change made in the browser. Ignore it. The instructor
      regenerates the notebooks after merging.\\[3pt]
    \textbf{Trope lint} &
      It never fails. Ignore it.\\[3pt]
    \textit{This branch has conflicts that must be resolved} &
      Do not try to fix it. The instructor resolves conflicts. You can still
      review the change.\\
    \bottomrule
  \end{tabular}\par}

  \handcapture{pr-checks.png}{Screenshot of the bottom of a pull request's
    Conversation tab, showing its checks and the merge box.}{The bottom of
    the \textbf{Conversation} tab: the checks, and the merge box. Do not
    click \textbf{Merge pull request}.}
\end{frame}

\begin{frame}{3 · Read the change}
  Click \textbf{Files changed}. Removed lines are red ①, and added lines are
  green ② (Figure~\ref{fig:diff}). Check five things:

  \begin{enumerate}
    \item \textbf{The right file.} Only the chapter's \texttt{.qmd} file
      changes. A file in \texttt{notebooks/} or \texttt{images/} must not
      change.
    \item \textbf{One change.} It changes what the description says, and
      nothing else.
    \item \textbf{Correct.} Open the same section of the book in a second
      tab, at \url{https://cuinfoscience.github.io/Web-Data-Science-Book/}.
      Are the facts right? Does each new link open the right page? Does new
      code use the chapter's names, such as \texttt{HEADERS}?
    \item \textbf{Clear.} Can a classmate who missed class understand it?
    \item \textbf{The Markdown.} The checks do not see these mistakes, so
      read for them:
      \begin{itemize}
        \item a link is \breakcode{[text](https://...)}, with no space between
          \texttt{]} and \texttt{(} ③;
        \item a heading (\texttt{\#\#}) or a callout (\texttt{:::}) needs an
          empty line above it ④. A heading straight after a callout's
          \texttt{:::} line is the callout's title, and needs none;
        \item a code block opens and closes with three backticks.
      \end{itemize}
  \end{enumerate}

  \activityfigure[0.86\linewidth]{handouts/week-07/img/diff-markdown_annotated.pdf}%
    {Screenshot of a change to chapter 4, line by line. Marker 1 is on the
     removed line, in red. Marker 2 is on the added line, in green. Marker 3
     is beside a Markdown link. Marker 4 is on the empty line above the
     callout.}%
    {One change, line by line (here week 4's example, in the web editor's
     Preview; \textbf{Files changed} uses the same colors). The removed line
     ① is red, and the added line ② is green. The new line writes a link as
     \breakcode{[text](https://...)} ③. The callout keeps the empty line above
     its \texttt{:::} line ④.}%
  \label{fig:diff}
\end{frame}

\begin{frame}{4 · Comment on a line}
  \begin{enumerate}
    \item Move the pointer over a line. A blue \textbf{+} shows beside it.
      Click it.
    \item Write one point: what is wrong, and what to change. To propose
      exact words, click \textbf{Add a suggestion} (the $\pm$ button), and
      edit the line in the box.
    \item Click \textbf{Start a review}. For your next comments, click
      \textbf{Add review comment}. Nobody sees them until you submit the
      review.
  \end{enumerate}
  \handcapture{pr-line-comment.png}{Screenshot of the Files changed tab with
    the comment box open on one line.}{A comment on one line of
    \textbf{Files changed}, before \textbf{Start a review}.}
\end{frame}

\begin{frame}{5 · Submit the review}
  \begin{columns}[T]
    \begin{column}{0.52\textwidth}
      \begin{enumerate}
        \item Click \textbf{Review changes}, at the top right of
          \textbf{Files changed}.
        \item Write a summary in the shape on the right.
        \item Select the option that matches your verdict:
          \begin{itemize}
            \item \textbf{Approve} for \textit{merge};
            \item \textbf{Request changes} for \textit{needs changes};
            \item \textbf{Comment} for \textit{close}, or if you are not sure.
          \end{itemize}
        \item Click \textbf{Submit review}.
      \end{enumerate}
    \end{column}
    \begin{column}{0.44\textwidth}
      \begin{block}{Your summary}
        \small\ttfamily
        Verdict: merge / needs changes / close\\
        Must fix: at most two points\\
        Nice to have: optional\\
        Reviewed with @partner
      \end{block}
    \end{column}
  \end{columns}

  Give the verdict \textit{close} for a duplicate (another pull request makes
  the same change) or for a change that the book does not need. Say why, and
  thank the author.

  A good review is kind and specific. It names the line, it keeps must-fix
  points apart from nice-to-have points, and it ends with a verdict.

  \handcapture{pr-review-changes.png}{Screenshot of the Review changes panel
    with its summary box and three options.}{\textbf{Review changes}: the
    summary, then \textbf{Comment}, \textbf{Approve}, or \textbf{Request
    changes}.}

  \medskip
  \textbf{Next.} Review the next pull request on your board. A pull request
  with a review already can use a second one.
\end{frame}
\end{document}
